\documentclass[11pt]{article}
\usepackage[margin=25mm]{geometry}
\usepackage{fontspec}
\setmainfont{Latin Modern Roman}
\usepackage{amsmath,amssymb,mathtools}
\usepackage{unicode-math}
\setmathfont{Latin Modern Math}
\usepackage{xurl,hyperref}
\hypersetup{colorlinks=true,urlcolor=blue}
\setcounter{secnumdepth}{0}
\setlength{\parindent}{0pt}
\setlength{\parskip}{5pt}
\setlength{\emergencystretch}{3em}
\newcommand{\sref}[2]{\hyperlink{source:#1:#2}{[S#2]}}
\newcommand{\cataloguescope}[1]{\paragraph{Recorded target.}#1}
\newcommand{\R}{\mathbb R}
\newcommand{\C}{\mathbb C}
\newcommand{\N}{\mathbb N}
\newcommand{\Z}{\mathbb Z}
\newcommand{\Q}{\mathbb Q}
\newcommand{\F}{\mathbb F}
\newcommand{\eps}{\varepsilon}
\begin{document}
% UnsolvedMath ID 30000180; source catalogue code OWR-785-006
\setcounter{section}{30000179}
\section{Positive Maps and Multipartite Bound Entanglement}\label{30000180}
{\small UnsolvedMath ID 30000180 · Mathematical Physics}
\subsection{Definitions and mathematical statement}

\paragraph{Shared notation.}$\mathbb N=\{1,2,\ldots\}$; the standard number fields have their usual meanings. All additional conventions are stated locally.


\paragraph{Definitions and conventions.} $M_d(\mathbb C)$ is the matrix algebra, with Hilbert--Schmidt inner product $\langle F,G\rangle=\operatorname{Tr}(F^\dagger G)$. A map is positive if it sends positive semidefinite matrices to positive semidefinite matrices, and completely positive if every tensor amplification by an identity map is positive. The source uses Hermitian orthonormal bases $F^{(i)}_\mu$ and real coefficients $\lambda^{(i)}_\mu$. For the four-qubit lattice states, let $\sigma_0=I_2$ and $\sigma_1,\sigma_2,\sigma_3$ be Pauli matrices; set $|\psi_+\rangle=\tfrac12\sum_{a,b=1}^2|a,b\rangle_A|a,b\rangle_B$ and $|\psi_{\alpha\beta}\rangle=(I_4\otimes\sigma_\alpha\otimes\sigma_\beta)|\psi_+\rangle$. Here $\varnothing\ne I\subseteq\{0,1,2,3\}^2$ and $|I|$ is its cardinality. PPT means positivity after partial transposition across the $A|B$ partition.

\paragraph{Source scope and qualifications.} The supplied positivity preamble omits the coefficient hypotheses on Lambda\_1 that are needed for the source's proposition. An arbitrary Lambda\_1 with large negative coefficients does not satisfy the stated conclusion. The primary source requires all coefficients of Lambda\_1 to be positive and sufficiently large relative to the single negative coefficient of Lambda\_2. This missing premise is explicitly restored in the source clarification.

\paragraph{Primary-source hypothesis.} In the report's positivity proposition, all $\lambda^{(1)}_\mu$ are positive and are at least the absolute value of the single negative coefficient of $\Lambda_2$; every other coefficient of $\Lambda_2$ has the same lower bound. The request is to generalize this family, for example by relaxing the orthonormal-basis identity, and to study the entanglement-detection properties of the states $\rho_I$.

\paragraph{Catalogue mathematical target.}
\paragraph{Catalogue wording retained for traceability.}
For $i=1,2$, let $$\Lambda_i[X]=\sum_{\mu=0}^{d_i^2-1}\lambda_\mu^{(i)}F_\mu^{(i)}XF_\mu^{(i)},$$ where the $F_\mu^{(i)}$ form Hermitian, Hilbert–Schmidt orthonormal bases of $M_{d_i}(\mathbb{C})$. If exactly one coefficient of $\Lambda_2$ is negative and every positive coefficient is at least its absolute value, then $\Lambda=\Lambda_1\otimes\operatorname{id}_{d_2}+\operatorname{id}_{d_1}\otimes\Lambda_2$ is positive. Can this family be generalized by relaxing the orthonormal-basis hypothesis? For $$\rho_I=\frac1{|I|}\sum_{(\alpha,\beta)\in I}|\psi_{\alpha\beta}\rangle\langle\psi_{\alpha\beta}|,$$ determine the properties of these states and, in particular, whether their entanglement can be detected by other criteria.

\paragraph{Saved literature context.} The supplied catalogue's assessment dated 2026-08-21 records: The report identifies open generalizations of the positive-map family and questions about the PPT bound-entangled states rho-tilde-I. The checked literature gives important criteria and constructions but no complete resolution of these exact questions. The assessment is retained as dated provenance; it is not a new comprehensive literature or open-status certification.
\paragraph{Scope and source qualification.}The supplied positivity preamble omits the coefficient hypotheses on Lambda_1 that are needed for the source's proposition. An arbitrary Lambda_1 with large negative coefficients does not satisfy the stated conclusion. The primary source requires all coefficients of Lambda_1 to be positive and sufficiently large relative to the single negative coefficient of Lambda_2. This missing premise is explicitly restored in the source clarification.
\subsection{Short English statement}
Extend the report's family of positive maps and determine how the associated Bell-diagonal states can exhibit detectable bound entanglement.
\subsection{Research attempt}
\paragraph{Model, date and investigated scope.}
OpenAI GPT-6 Astra (configured model \texttt{gpt-6-astra}), ultra reasoning; generated 11 October 2026. Tests redundant tight Hermitian frames and proves a simple Bell-overlap witness for lattice states of cardinality below four.

\paragraph{Sources and literature check.}
The complete supplied definition, including source qualifications and the dated literature assessment, was reviewed. The calculation below is an elementary restricted attempt motivated by [S2]; it does not claim novelty or a fresh literature-status determination.

\paragraph{Context and investigated scope.} The complete supplied definition, including source qualifications and the dated literature assessment, was reviewed. The calculation below is an elementary restricted attempt motivated by [S2]; it does not claim novelty or a fresh literature-status determination.

\paragraph{A controlled relaxation of orthonormality.} Let $F_1,\ldots,F_{d^2}$ be a Hermitian Hilbert--Schmidt orthonormal basis and let $O$ be a real $N\times d^2$ matrix with $O^TO=I$, where $N\geq d^2$. Define $G_a=\sum_\mu O_{a\mu}F_\mu$. These are a possibly redundant, nonorthogonal Hermitian family and satisfy
\[
 \sum_aG_aXG_a=\sum_{\mu,\nu}(O^TO)_{\mu\nu}F_\mu XF_\nu
 =\sum_\mu F_\mu XF_\mu=\operatorname{Tr}(X)I.
\]
The last identity can be checked with diagonal matrix units and the symmetric/antisymmetric Hermitian combinations of off-diagonal units, and is preserved by real orthogonal basis changes. Duplicating each $F_\mu/\sqrt2$ gives a concrete redundant example. Thus any proof using only this unweighted completeness identity can retain that step for these frames. It cannot automatically retain a theorem with independently prescribed weighted coefficients: mixing introduces cross terms unless the coefficient matrix is transformed too.

\paragraph{A lattice-state entanglement test.} Each $|\psi_{\alpha\beta}\rangle$ is maximally entangled across $\mathbb C^4|\mathbb C^4$. For normalized product vectors $u\otimes v$, Cauchy--Schwarz gives $|\langle\psi_{\alpha\beta}|u\otimes v\rangle|^2\leq1/4$. Convexity extends this to separable states. Orthogonality of Pauli matrices makes the sixteen Bell vectors orthonormal, hence
\[
 \langle\psi_{\alpha\beta}|\rho_I|\psi_{\alpha\beta}\rangle=1/|I|\quad((\alpha,\beta)\in I).
\]
Consequently $|I|<4$ is an elementary entanglement certificate. At $|I|=4$ it becomes inconclusive, not a separability proof; at $|I|=16$ the state is $I_{16}/16$ and separable.

\paragraph{Unresolved parts.} The frame identity does not prove positivity for arbitrary nonorthogonal weighted families. The source's restored positive lower bounds on the first map remain essential premises and have not been dropped. The fidelity witness does not classify the PPT lattice states or detect their bound entanglement in general. These two gaps correspond to the two original requests.

\paragraph{Verification and adversarial audit.}
Checked the frame identity, necessity of transforming weighted coefficients, separable overlap ceiling, cardinalities 1, 4 and 16, and the restored first-map coefficient assumptions. Model self-audit only; no independent human proof checking. Prompts A and B were applied to this canonical record and its qualified source scope. Checked the frame identity, necessity of transforming weighted coefficients, separable overlap ceiling, cardinalities 1, 4 and 16, and the restored first-map coefficient assumptions. Retained unresolved status and zero original-target progress; Weighted nonorthogonal-map positivity and detection/classification of the PPT lattice states remain unproved.

Coordinator checked the tight-frame completeness identity, weighted-family limitation and 1/4 product-overlap bound for dimension-four maximally entangled states. This additional check was model-assisted, not independent human verification.

\paragraph{Tools, process and contributions.}
Python 3 was used for file preparation and structural checks; no numerical experiment or theorem prover was used. Contribution prompt A was followed by prompt B. GPT-6 Astra with ultra reasoning drafted the research and performed its stated self-audit; any additional coordinator contribution is identified in the verification and audit. A GPT-6-based Codex coordinator assembled the manuscript and checked repository integration; its exact variant and reasoning configuration were not exposed. No human mathematical contribution, independent human proof check or formal proof certification is supplied. Complete derivations and the consulted source references are included here; private reasoning transcripts are not published.

\paragraph{Final status and first remaining gap.}
\textbf{UNRESOLVED}. Weighted nonorthogonal-map positivity and detection/classification of the PPT lattice states remain unproved.

\section{Completion Estimate}
\noindent\textbf{COMPLETION: 0\%}
The calculation is an auxiliary or restricted attempt and establishes no newly verified portion of the full source target. Weighted nonorthogonal-map positivity and detection/classification of the PPT lattice states remain unproved.

\subsection{Sources}
\begin{itemize}
\item[\hypertarget{source:30000180:1}{S1}] ULAM.AI and the UnsolvedMath Contributors, UnsolvedMath, record 30000180 (OWR-785-006). Catalogue curation under CC BY 4.0; definitions and English exposition added locally.
\url{https://huggingface.co/datasets/ulamai/UnsolvedMath}.
\item[\hypertarget{source:30000180:2}{S2}] Entanglement and Decoherence: Mathematics and Physics of Quantum Information and Computation (2005), Oberwolfach Reports, DOI 10.4171/owr/2005/04, pp. 185--255; official report source and mathematical context.
\url{https://doi.org/10.4171/owr/2005/04}.
\item[\hypertarget{source:30000180:3}{S3}] Official EMS Press PDF of the same Oberwolfach report; PDF pages 48--50, printed pages 232--234.
\url{https://ems.press/content/serial-article-files/45978}.
\item[\hypertarget{source:30000180:4}{S4}] Horodecki et al., multipartite separability and positive maps; reference in the saved dated catalogue literature review.
\url{https://arxiv.org/abs/quant-ph/0006071}.
\end{itemize}
\label{end:30000180}
\end{document}
